\subsection{DEFINITION OF QUANTIFIERS}

The only logical signs which played a role in § 3 are \(\neg\) and \(\vee\). The rules we shall now state are essentially concerned with the use of the logical signs \(\tau\) and \(\square\).

If R is an assembly and x a letter, the assembly \((\tau_{x}(R)|x)R\) is denoted by ``there exists x such that \(R\)'', or by \((\exists x)R\) . The assembly

\[
\text {   ``not   } ((\exists x) \text {   not   } R)
\]

is denoted by ``for all x, R'', or by ``given any x, R'', or by \((\forall x)R\) . The abbreviating symbols \(\exists\) and \(\forall\) are called respectively the existential quantifier and the universal quantifier. The letter x does not appear in the assembly denoted by \(\tau_{x}(R)\) and therefore does not appear in the assemblies denoted by \((\exists x)R\) and \((\forall x)R\) .

CS8. Let R be an assembly and let x and \(x'\) be letters. If \(x'\) does not appear in R, then \((\exists x)R\) and \((\forall x)R\) are respectively identical with \((\exists x')R'\) and \((\forall x')R'\) , where \(R'\) is \((x'|x)R\) .

For \((\tau_{\mathbf{x}}(\mathbf{R})|\mathbf{x})\mathbf{R}\) is identical with \((\tau_{\mathbf{x}}(\mathbf{R})|\mathbf{x}')\mathbf{R}'\) by CS1 (§1, no. 2), and \(\tau_{\mathbf{x}}(\mathbf{R})\) is identical with \(\tau_{\mathbf{x}'}(\mathbf{R}')\) by CS3 (§1, no. 2). Hence \((\exists \mathbf{x})\mathbf{R}\) is identical with \((\exists \mathbf{x}')\mathbf{R}'\) . It follows that \((\forall \mathbf{x})\mathbf{R}\) is identical with \((\forall \mathbf{x}')\mathbf{R}'\) .

CS9. Let R and U be assemblies and let x and y be distinct letters. If x does not appear in U, then \((U|y)(\exists x)R\) and \((U|y)(\forall x)R\) are respectively identical with \((\exists x)R'\) and \((\forall x)R'\) , where \(R'\) is \((U|y)R\) .

For, by CS2 (§ 1, no. 2), \((\pmb{U}|\pmb{y})(\tau_{\pmb{x}}(\pmb{R})|\pmb{x})\pmb{R}\) is identical with

\[
(T | \boldsymbol {x}) (U | \boldsymbol {y}) R,
\]

where \(T\) is \((U|y)\tau_x(R)\) , that is \(\tau_x(R')\) , by CS4 (§1, no. 2). Hence \((U|y)(\exists x)R\) is identical with \((\exists x)R'\) , and consequently \((U|y)(\forall x)R\) is identical with \((\forall x)R'\) .

CF11. If R is a relation in a theory C and if x is a letter, then \((\exists x)R\) and \((\forall x)R\) are relations in C.

This follows immediately from CF3, CF8, and CF2 (§ 1, no. 4).

Intuitively, let us consider R as expressing a property of the object denoted by x. From the intuitive meaning of the term \(\tau_{x}(R)\) , the assertion \((\exists x)R\) means that there is an object which has the property R. The assertion ``not \((\exists x)(\text{not } R)\) '' means that there is no object which has the property ``not R'', and therefore that every object has the property R.

In a logical theory \(\mathcal{C}\) , if we have a theorem of the form \((\exists x)R\) , in which the letter \(x\) is not a constant of \(\mathcal{C}\) , this theorem can serve as a theorem of legitimation in the method of the auxiliary constant ( \(\S 3\) , no. 3), because it is identical with \((\tau_{x}(R)|x)R\) . Let \(\mathcal{C}'\) be the theory obtained by adjoining \(R\) to the axioms of \(\mathcal{C}\) . If we can prove a relation \(S\) , in which \(x\) does not appear, in the theory \(\mathcal{C}'\) , then \(S\) is a theorem in \(\mathcal{C}\) .

C26. Let \(\mathfrak{G}\) be a logical theory, let \(\pmb{R}\) be a relation in \(\mathfrak{G}\) , and let \(\pmb{x}\) be a letter. The relations \((\forall x)\pmb{R}\) and \((\tau_{x}(\mathrm{not}\pmb{R})|\pmb{x})\pmb{R}\) are then equivalent in \(\mathfrak{G}\) .

For \((\forall x)R\) is identical with ``not \((\tau_{x}(\text{not } R)|x)(\text{not } R)\) '' and therefore with ``not not \((\tau_{x}(\text{not } R)|x)R\) ''.

C27. If R is a theorem in a logical theory G in which the letter x is not a constant, then \((\forall x)R\) is a theorem in G.

For \((\tau_{\pmb{x}}(\text{not } \pmb{R})|\pmb{x})\pmb{R}\) is a theorem in \(\mathfrak{C}\) , by C3 (§2, no. 3).

On the other hand, if x is a constant of G, the truth of R does not imply that of \((\forall x)R\) . Intuitively, the fact that R is a true property of x, which is a definite object of G, clearly does not imply that R is a true property of every object.

C28. Let \(\mathfrak{G}\) be a logical theory, let \(\pmb{R}\) be a relation in \(\mathfrak{G}\) , and let \(\pmb{x}\) be a letter. Then the relations ``not \((\forall x)\pmb{R}\)'' and \((\exists x)\) (not \(\pmb{R}\) ) are equivalent in \(\mathfrak{G}\) .

For ``not \((\forall x)R\) '' is identical with ``not not \((\exists x)(\text{not } R)\) ''.

